% =====================================================================
% SEC_01 — Introdução e Escopo
% =====================================================================

\section*{1. Introdução e Escopo}
\addcontentsline{toc}{section}{1. Introdução e Escopo}

\nivel{0} Este manual tem duas finalidades complementares. A primeira é documentar,
de forma precisa e verificável, a arquitetura de Recuperação Aumentada por
Geração (RAG) tal como \emph{implementada} no ecossistema, referenciando os
módulos reais do código-fonte. A segunda é apresentar uma \textbf{proposta} de
evolução dessa arquitetura: a introdução de um estágio de diversificação
estruturada baseado na sequência matemática de Recamán, que atua sobre o
ranqueamento dos candidatos recuperados.

\begin{intuicao}
Pense no RAG como um assistente de pesquisa que, antes de responder, consulta
uma biblioteca. A qualidade da resposta depende de \emph{quais livros} ele
coloca sobre a mesa. O ponto central deste manual é simples: às vezes a mesa fica
cheia de livros muito parecidos entre si. A \emph{diversificação} tenta resolver
isso, garantindo que os livros representem ângulos diferentes do tema.
\end{intuicao}

\nivel{2} É fundamental registrar uma distinção metodológica que orienta todo o
documento: a arquitetura \emph{atual} é observável e auditável no checkout,
enquanto a arquitetura \emph{pós-Recamán} configura um desenho-alvo
\textbf{proposto}, ainda não implementado. Esta separação entre o que \emph{é} e
o que \emph{poderia ser} sustenta o rigor anti-overclaim do ecossistema
(SDD/TDD): nenhum resultado é apresentado como implementado sem ser observável no
código. Nenhuma alegação de certificação externa, de segurança absoluta ou de
disponibilidade universal de serviços é feita aqui.

\begin{pontochave}
Ao final desta leitura, você deve ser capaz de: (i) distinguir o estado atual da
proposta pós-Recamán; (ii) identificar o problema de pesquisa que motiva a
diversificação; (iii) compreender a intuição por trás do uso da sequência de
Recamán; e (iv) localizar as especificações técnicas e referências que sustentam
cada afirmação.
\end{pontochave}

\subsection*{1.1 Como ler este manual (trilha de profundidade)}

\nivel{0} Este manual foi escrito para ser lido por camadas. Cada seção é
marcada com um \textbf{nível de profundidade} (N0, N1 ou N2) que indica o grau de
formalização exigido:

\begin{center}
\begin{tabularx}{\textwidth}{llX}
\toprule
\textbf{Nível} & \textbf{Tipo de leitor} & \textbf{Em que consiste} \\
\midrule
N0 & Introdutório & Intuição, analogias e motivação; a "história" do problema. \\
N1 & Aplicado & Como os conceitos se materializam em componentes; leitura de
diagramas, tabelas e trechos de código. \\
N2 & Pós-graduação & Formalização matemática, equações, complexidade
assintótica e diálogo com a literatura (referências reais). \\
\bottomrule
\end{tabularx}
\end{center}

\noindent O leitor de nível 0 pode percorrer apenas os parágrafos marcados N0 e
as \emph{Caixas de intuição}. O leitor de nível PhD deve percorrer tudo, incluindo
as equações e as referências. Nenhuma informação é "escondida": a mesma ideia
aparece em várias profundidades, de forma didática, sem sacrificar o rigor da
publicação.

\subsection*{1.2 Estrutura do documento}

O manual está organizado da seguinte forma:

\begin{enumerate}[label=\arabic*.]
  \item \textbf{Introdução e escopo} --- esta seção.
  \item \textbf{Problema de pesquisa e impactos} --- formulação do problema que
        motiva o estudo e seus impactos na ciência e no ecossistema.
  \item \textbf{Fundamentação teórica} --- modelos de ranqueamento, RAG e MMR.
  \item \textbf{Estado atual da arquitetura RAG} --- componentes implementados.
  \item \textbf{Proposta pós-Recamán} --- arquitetura-alvo e diversificador.
  \item \textbf{Cálculos e especificações técnicas} --- equações e métricas.
  \item \textbf{Implementações reproduzíveis} --- pseudocódigo e trechos.
  \item \textbf{Mapas de dados e diagrama operacional} --- visualização.
  \item \textbf{Considerações finais e trabalho futuro}.
\end{enumerate}

\subsection*{1.3 Princípios de rigor}

\nivel{2} Seguindo o protocolo de rigor do ecossistema (SDD/TDD e anti-overclaim),
todo resultado aqui apresentado é ou (i) observável no código, ou (ii)
reproduzível a partir das especificações, ou (iii) claramente rotulado como
proposta. As referências bibliográficas correspondem a fontes reais validadas ---
DOIs de periódicos/conferências e preprints arXiv --- e nenhuma foi fabricada.
A didática nunca substitui a precisão: analogias são sempre acompanhadas da
formalização correspondente.
